\cleardoublepage

\newrefsection

\chapter{外文翻译}

\sectionnonum{讨论}

\par 本研究提供了证据，表明通过AI驱动的对话代理进行每日目标规划，相比脚本化的非AI聊天机器人对照组，在三个主要结果上产生了更大的益处。接受生成式AI聊天机器人辅导的参与者表现出生产力和执行意图的显著更大改善，以及拖延症的更大减少。在探索性结果中，AI组相对于对照组还导致了低唤醒积极情感的显著增加，以及高唤醒消极情感和心智游移的显著减少。然而，几个主要结果在两组间改善相似，包括目标进展、目标实现、行动和应对计划，以及工作联盟。无论参与者的基线特征如何，都观察到了这些结果。探索性调节分析表明，基线目标承诺、AI依赖、对AI的态度和对AI的信任并未调节条件对任何主要结果的影响，除了AI依赖调节了条件对目标实现的影响。然而，基线AI依赖显著调节了目标实现的变化，使得AI条件仅对低依赖性的参与者有益。下面讨论这些发现的含义。

\par 本发现对自我调节和目标追求的理论模型做出了贡献，强调结构化每日目标规划是一种有效策略，并表明适应性AI对话可能进一步加强关键自我调节机制。与自我调节框架一致，两组在干预周内都表现出计划过程（行动和应对计划）和目标相关结果（目标进展和目标实现）的可比改善。这种模式表明，将每周目标分解为可操作任务有助于缩小意图-行为差距（Carver \& Scheier, 1998; Sheeran \& Webb, 2016），支持改善的计划和目标特定结果，无论交付模式如何。然而，AI组在执行意图的制定方面提供了额外益处，操作化为采取具体步骤实现目标，相比于非AI聊天机器人对照组。尽管两组参与者都参与了时间盒和if-then计划，并在行动和应对计划结果上显示出可比的收益，AI组在支持行为跟进方面更有效。与执行意图理论一致，这种益处可能反映了情境-行动联结的加强（Cohen \& Hicks, 2017; Gollwitzer, 1999），得到AI聊天机器人情境意识的支持，推荐情境相关行动并加强对目标导向行为的责任感（Jones et al., 2024）。更一般地说，这种效应可能反映与AI聊天机器人更丰富质的参与，这可能促发面向未来的思维和持续行动（Dechant et al., 2025）。

\par 除了支持执行意图外，AI支持还加强了更广泛的自我调节参与，反映在生产力的更大改善上，相比于对照组。这些收益可能反映了在干预周内感知到的成就和任务完成的增加。当前干预中使用的AI聊天机器人支持每日目标规划，可能减少了与常规规划任务相关的认知负荷，使个体能够将认知资源分配给更复杂的生产力任务（Iqbal et al., 2025; Marois et al., 2025）。同时，AI组的参与者在拖延上表现出更大的减少，表明相对于对照组，在启动意图任务上的延迟更少。这些发现表明AI辅导可以帮助缓解自我调节失败，其特征是为短期诱惑而延迟意图任务（Steel, 2007; Sirois \& Pychyl, 2013）。作为由大语言模型驱动的系统，AI聊天机器人可以将较大任务分解为可管理的组成部分并提供可操作的下一步，从而促进任务启动并减少拖延（Bhattacharjee et al., 2023）。

\par 虽然拖延反映了意图行动过程中的延迟（Steel, 2007; Lay \& Silverman, 1996），探索性分析表明AI组的参与者经历了心智游移的更大减少，相比于对照组。心智游移反映了任务无关思维对任务聚焦注意力的中断（Mrazek et al., 2013; Seli et al., 2015），表明AI聊天机器人的结构化指导可能支持了目标追求期间的持续注意力投入。一个合理的机制是，AI的决策支持，如预测相关下一步和提供具体建议，减少了认知负荷并减弱了注意力漂移（Gandhi et al., 2023; Schmidhuber et al., 2021）。重要的是，心智游移的减少发生在没有相应的一般认知失败变化的情况下，表明这些注意力益处可能独立于更广泛的认知失败而发生（McVay \& Kane, 2010）。总体而言，这些发现表明对话AI可能充当促进持续注意力的认知支持系统。这些提出的机制仍然是推测性的，因为认知负荷和注意力过程在当前研究中未被直接测量，应在未来研究中更直接地检验。

\par 除了行为和注意力结果外，AI辅导聊天机器人还产生了情感益处，包括高唤醒消极情感（即愤怒、敌意、易怒）的显著减少和低唤醒积极情感（即平静、满足、放松）的增加，两者都发生在相似的唤醒维度上。减少的高唤醒消极情感可能反映了目标追求期间由不确定性产生的摩擦和挫折的降低，得到AI聊天机器人个性化回应和量身定制验证的支持（Schmidhuber et al., 2021; Hu et al., 2025）。相反，低唤醒积极情感的增加可能反映了更大的情感轻松，由不施加紧迫感的鼓励和非评判性支持所培养（Inkster et al., 2018; Terblanche et al., 2022）。这些结果与先前研究一致，表明易于获取且响应迅速的AI系统可以为幸福感和情绪自我调节做出贡献（Alabed et al., 2024; Passmore \& Tee, 2023; Sun et al., 2025）。这些情感改善表明AI目标规划干预可能将情绪状态转向更大的平静，进而可能通过促进专注和减少压力来支持对目标追求的持续参与。

\par 本研究的另一个关键发现是，尽管AI聊天机器人被设计为传达共情，工作联盟从前测到后测保持稳定，随时间没有增加，变化也没有组间差异。工作联盟捕捉了与聊天机器人在三个维度上的协作伙伴关系——目标、任务和纽带——分别反映了对推进目标的信任、对目标导向协作步骤的协议以及关系感和支持（Bordin, 1979; Goldberg et al., 2022）。与先前工作一致，自由文本AI聊天机器人与纽带的适度增加相关，但与基于点击的预定义响应聊天机器人相比，在整体工作联盟上没有明显优势（Mai et al., 2022）。更强伙伴关系的发展可能需要比本研究5天持续时间更长的互动期，因为证据有限表明用户在如此短暂的时间范围内与AI对话代理建立纽带（Beatty et al., 2022）。尽管如此，先前研究发现，即使与AI聊天机器人的工作联盟未能发展，大多数参与者仍能改善他们的自我韧性（Ellis-Brush, 2021）。本研究观察到某些自我调节结果的益处，表明个体仍可在没有强烈工作联盟的情况下接受并受益于AI指导。

\par 观察到的目标追求和探索性幸福感结果的收益突出了基于AI聊天机器人的目标规划作为可扩展工具的实践潜力，用于在不同情境下支持生产力和日常幸福感。例如，此类AI元素可以整合到消息平台或现有的生产力工具中，如待办事项或日历应用程序，以支持学生的自我调节行为。支持这种方法的可行性，将半生成式聊天机器人整合到待办事项应用程序中导致在提供心理支持和减少压力方面的显著改善（Lee et al., 2025）。类似的AI支持干预也可能适用于组织情境，其中AI工具可以提高工作生产力和绩效（Al Naqbi et al., 2024）。除了生产力增强外，关于幸福感的探索性发现对AI辅助辅导在情感支持情境中的使用具有实际意义，扩展了先前关于专注于改善幸福感和个人成长的AI干预的研究（Blyler \& Seligman, 2024; Shin, 2025）。然而，当前干预是简短的，且目标规划任务相对简单。更复杂的辅导领域，如开放式咨询或深度职业辅导，可能需要进一步的发展和伦理保障。

\par 尽管AI聊天机器人交付的辅导具有成本效益、可扩展性和持续可用性等优势，但关于不准确性和保密性伦理关切的限制仍然存在（Ho et al., 2025; Passmore \& Tee, 2023; Terblanche, 2022）。考虑到这些限制和人类教练的优势，包括细致的情感理解，AI-人类混合方法可能是有前景的。在此类模式中，可在会话之间使用的AI基础辅导可以补充更深入的人工辅导会话（Bachkirova \& Kemp, 2025; Bruning \& Boak, 2025）。早期尝试正在探索与人工教练一起的AI增强聊天机器人（Arakawa \& Yakura, 2024; Bojic et al., 2025），表明对利用AI和人类教练互补优势的混合辅导模式的兴趣和潜力不断增长。

\par 在解释本发现时应考虑几个局限性。样本由大学生组成，干预集中于为期一周的学术或个人目标。这些可能限制了向更广泛人群的推广，具有不同的情境需求或目标结构（Peterson \& Merunka, 2014; Shadish et al., 2002）。这个相对同质的群体可能特别容易受到拖延等自我调节挑战的影响（Fentaw et al., 2022; Özer, 2011）。作为数字原住民，他们也可能比其他人群更精通技术，更熟悉基于AI聊天机器人的干预（Jones et al., 2010; Singh et al., 2023）。进一步的研究应检验基于AI的目标规划干预在更多样化的具有自我调节挑战的群体中的有效性，如具有高工作需求的职场专业人士（Bakker \& de Vries, 2021），以及包括注意缺陷多动障碍个体的临床人群，他们可能从AI驱动的情绪调节策略中受益（Faraone et al., 2019; Thakkar et al., 2024）。

\par 此外，干预是简短的，仅跨越5天的积极使用，这限制了关于其长期有效性的结论。目前尚不清楚观察到的收益在AI支持撤回后是否会持续，或随着新奇感消退而减弱。更长的干预设计——跨越数周或整个学期——结合后续评估，需要评估随时间的用户参与变化、潜在的新奇效应，以及持续互动在多大程度上支持习惯形成和关系过程（Sadeghi et al., 2022）。当前研究的结果依赖于自我报告测量，这可能容易受到社会期望偏差的影响（Paulhus \& Vazire, 2007）。尽管组分配大体成功，但在行动计划、目标承诺和对AI的信任方面观察到预先存在的基线差异，AI组的参与者表现出行动计划的轻微初始优势。由于这些变量未进行统计控制，它们对结果的影响无法完全排除。然而，探索性调节分析发现，AI态度和AI信任并未显著调节条件对所有主要结果的影响，表明AI每日目标规划的收益在很大程度上对这些基线差异具有稳健性。

\par 除了基线考虑外，未来研究将受益于纳入目标追求的更客观指标，如学术表现或任务完成测量，以补充自我报告结果并加强因果推断（Smyth et al., 2023）。

\par 未来研究可以通过将AI目标规划干预从独立聊天转移到用户现有的数字生态系统中，并进一步探索关系设计元素来扩展当前工作。将AI教练与任务管理系统或日历整合可以启用更特定情境的提示，例如通过识别可用时间窗口（Myers et al., 2007; Rajamoorthy et al., 2025）。这种整合还将允许监控客观指标，这可以支持适应性学习并实现对个体用户的辅导策略个性化。进一步的工作可以检验用户定制选项，包括激励性与分析性沟通风格的有效性。虽然当前干预采用了主要任务聚焦的设计，未来研究可以探索具有更大社会存在感和情感连接的AI干预，用于更情感导向的干预。然而，高度拟人化设计特征的使用可能引起伦理关切，突出持续评估信任和隐私问题的重要性（Osika, 2025; Salles et al., 2020）。

\par 总之，本研究突出了AI驱动的每日目标规划作为提高生产力、减少拖延和改善幸福感的有效工具的潜力。尽管干预是简短的，但发现支持AI辅助辅导可能有助于解决自我调节挑战。具体而言，AI聊天机器人提供了个性化指导，帮助学生将目标转化为行动，贯彻意图，并比基于规则的非AI聊天机器人更有效地管理分心和情绪。这些结果通过突出其作为增强某些目标追求结果工具的潜力，为日益增长的AI辅导文献做出了贡献。未来研究应关注AI干预的长期有效性、它们整合到现有数字工具中，以及利用AI和人类指导双方优势的混合方法的发展，同时保持对伦理关切的警觉。

\begingroup
    \linespreadsingle{}
    \printbibliography[title={外文翻译参考文献}]
\endgroup
